\section{Conclusion}\label{sec:conclusion}

\begin{callout}
\textbf{One endpoint, your machines only.} Ymir treats a handful of personal machines as a single
private model cloud: one command to start, a link to grow, no ports or keys to expose, and answers
that survive a node dying halfway through.
\end{callout}

Its core bet is that identity-keyed peer-to-peer transport can replace the proxy, TLS, API-key and
VPN layers that self-hosters normally bolt on, and that the generic pieces (discovery, inference
protocol, sync, engines) belong in shared Weirdnet crates so any app can use the cloud. The
single-box path, multi-node routing, context sync and two-level failover are built; pool-wide
placement, Ray-based splitting and cert-only service nodes are the documented road ahead.

Ymir is developed by Michael Christenson II and published by \WB{} under the BSD-3-Clause license.
Copyright 2026 \WB.

\begin{itemize}
  \item Weird Brains: \url{https://weirdbrains.com}
  \item Ymir on GitHub: \url{https://github.com/WeirdBrains/ymir}
  \item Ymir design doc: \url{https://github.com/WeirdBrains/ymir/blob/main/docs/DESIGN.md}
  \item Bragi on GitHub: \url{https://github.com/WeirdBrains/bragi}
  \item Weirdnet: \url{https://net.weirdbrains.com}
\end{itemize}
